\documentclass[12pt, a4paper]{article}

% --- Packages ---
\usepackage[utf8]{inputenc}
\usepackage[T1]{fontenc}
\usepackage{geometry}
\geometry{top=2.5cm, bottom=2.5cm, left=3cm, right=2.5cm}
\usepackage{amsmath, amssymb}
\usepackage{graphicx}
\usepackage{hyperref}
\usepackage{cite}
\usepackage{setspace}
\usepackage{booktabs}
\usepackage{array}
\usepackage{multirow}
\usepackage{CJKutf8}
\onehalfspacing

\hypersetup{
    colorlinks=true,
    linkcolor=blue,
    filecolor=magenta,
    urlcolor=blue,
    citecolor=blue,
}

% --- Document Identity ---
\title{\textbf{Visual Recognition using Deep Learning}\\ \large Homework 4: Image Restoration}
\author{
    Muhammad Rayhan Athaillah \begin{CJK*}{UTF8}{bsmi}(賴瑞涵)\end{CJK*}\\
    \small Student ID: 313540001\\
    \small National Yang Ming Chiao Tung University
}
\date{\today}

\begin{document}

\maketitle

\vspace{0.5cm}
\noindent \textbf{GitHub Repository:} \href{https://github.com/RayhanHaqi/visual\_recognition-hw4}{https://github.com/RayhanHaqi/visual\_recognition-hw4}

% ==========================================
% 1. INTRODUCTION
% ==========================================
\section{Introduction}
This report presents the technical implementation for the Image Restoration problem (Homework 4) within the Visual Recognition using Deep Learning course. The objective is to train a single PromptIR model \cite{potlapalli2023} capable of restoring images degraded by rain streaks and snow overlay, evaluated by PSNR on the CodaBench competition platform.

\subsection{Task and Dataset}
The dataset consists of 3,200 paired training images (1,600 per degradation type: rain and snow) and 320 validation images (160 per type), each at $256 \times 256$ resolution. Every degradation type has corresponding clean ground-truth images. The test set contains 100 degraded images (50 rain, 50 snow) without clean labels, submitted as a single \texttt{pred.npz} inside a \texttt{.zip} archive to CodaBench.

\subsection{Technical Constraints}
The assignment enforces the following constraints:
\begin{itemize}
    \item The model must be based on \textbf{PromptIR} \cite{potlapalli2023}; architecture modifications to PromptIR components/modules are permitted, but the core PromptIR model must remain.
    \item \textbf{No external data} beyond the provided dataset is allowed.
    \item \textbf{No pretrained weights} — the model must be trained from scratch.
    \item A \textbf{single model} must handle both rain and snow degradation types.
    \item All source code must be submitted as \texttt{.py} files alongside a PDF report in English.
\end{itemize}

% ==========================================
% 2. LITERATURE REVIEW
% ==========================================
\section{Literature Review}

\subsection{PromptIR: Prompting for All-in-One Blind Image Restoration}
PromptIR, proposed by Potlapalli et al. \cite{potlapalli2023}, introduces a prompt-based learning approach for all-in-one image restoration. Its key design contributions include:

\begin{itemize}
    \item \textbf{Prompt Module:} Learnable prompt components encode degradation-specific information through cross-attention with input features. The prompts are produced by a lightweight Prompt Generation Module (PGM) that interacts with a set of frozen prompt tokens.
    \item \textbf{Transformer-Based Restoration Backbone:} A U-Net style encoder-decoder architecture with transformer blocks that are dynamically guided by the generated prompts. The prompts modulate the feature processing at multiple scales, enabling the network to adapt its behavior to different degradation types without type-specific branches.
    \item \textbf{All-in-One Capability:} By learning a shared prompt space, PromptIR can restore images from multiple degradation types (denoising, deraining, dehazing) with a single unified model, generalizing to unseen degradation combinations.
\end{itemize}

The model contains approximately 35.6 million parameters and is trained with L1 loss, making it competitive with task-specific methods while maintaining the flexibility of a single restoration network.

\subsection{Related Work: All-in-One Image Restoration}
Several recent methods have addressed the all-in-one restoration challenge. TransWeather \cite{valanarasu2022} uses a single encoder-decoder transformer with learnable weather-type embeddings to handle rain, fog, and snow. Restormer \cite{zamir2022} introduced efficient multi-head transposed attention and gated feed-forward networks for high-resolution restoration tasks, achieving state-of-the-art on individual tasks but requiring task-specific training. AirNet \cite{li2022} learns degradation representations to guide a shared restoration network. However, these architectures are not permitted for this assignment, which specifically mandates PromptIR as the model backbone.

\subsection{Test-Time Augmentation and Checkpoint Averaging}
Test-time augmentation (TTA) \cite{ashukha2020} improves prediction robustness by averaging model outputs across multiple geometric transformations of the input. For restoration tasks, 8-way TTA combining identity, horizontal/vertical flips, and transpositions is a common technique. Checkpoint averaging (model soup) \cite{horton2021} averages weights from multiple nearby checkpoints to find a flatter minimum in the loss landscape, often producing better generalization than any single checkpoint.

% ==========================================
% 3. METHODOLOGY
% ==========================================
\section{Method}

\subsection{Data Preprocessing and Augmentation}
The dataset is loaded using \texttt{torchvision.datasets.ImageFolder}-style directory structure. During training, images are cropped to a patch of size $P \times P$ (with ground-truth classification to ensure proper loss computation), and random D4-group augmentation is applied: one of 8 transformations (identity, horizontal flip, vertical flip, both flips, and their transposed variants) is selected randomly. The augmentation probability includes identity mode to avoid over-augmentation. Images are cropped to multiples of 16 before being tensorized and normalized to $[0, 1]$. Validation and test sets are loaded as full $256 \times 256$ images with no augmentation.

\subsection{Model Architecture: PromptIR}
The PromptIR model \cite{potlapalli2023} is used as the core restoration network. It consists of:

\begin{itemize}
    \item \textbf{Prompt Generation Module (PGM):} Takes a set of learnable prompt components and a query derived from input features. Applies cross-attention to produce degradation-specific prompts.
    \item \textbf{Encoder-Decoder Transformer Blocks:} Each block contains multi-head self-attention and feed-forward layers, modulated by the degradation prompt via spatial feature transformation.
    \item \textbf{Skip Connections:} Multi-scale encoder features are fused with decoder features via prompt-guided feature refinement.
    \item \textbf{Output Head:} A convolutional layer maps decoder features back to a 3-channel restored image.
\end{itemize}

The total model size is 35.6 million trainable parameters, well-suited for $256 \times 256$ image restoration on a single RTX 4090 GPU with batch size 1.

\subsection{Training Pipeline}
The default training configuration uses:
\begin{itemize}
    \item \textbf{Loss:} L1 loss between restored and clean images (later upgraded to $L1 + \alpha\cdot\text{MSE}$).
    \item \textbf{Optimizer:} AdamW with learning rate $2 \times 10^{-4}$.
    \item \textbf{LR Schedule:} LinearWarmupCosineAnnealingLR with 15 epochs of linear warmup followed by cosine annealing to near-zero over 150 epochs.
    \item \textbf{Batch Size:} 1 (full $256 \times 256$ images with no gradient accumulation).
    \item \textbf{Checkpointing:} Saved every 5 epochs with Top-K monitoring of \texttt{val\_psnr} (later switch from \texttt{val\_loss}).
    \item \textbf{Inference:} Full-image feed-forward with optional 8-way geometric TTA and reflect-padding to the nearest multiple of 16.
    \item \textbf{Output:} Rounded uint8 conversion from $[0,1]$ float tensors using \texttt{np.rint} instead of flooring.
\end{itemize}

\subsection{Modifications and Improvements}

\subsubsection{PSNR-Aligned Loss Functions}
The original PromptIR uses pure L1 loss. Since the evaluation metric is PSNR (which is MSE-derived), a hybrid loss was implemented: $\mathcal{L}_{\text{L1+MSE}} = \mathcal{L}_{\text{L1}} + \lambda \cdot \mathcal{L}_{\text{MSE}}$, with configurable weight $\lambda$. A Charbonnier loss (smooth L1 variant) was also added as an alternative. The system supports \texttt{--loss\_type l1|l1\_mse|charbonnier} and \texttt{--mse\_weight} flags.

\subsubsection{Task-Aware Conditioning}
A lightweight wrapper (\texttt{TaskConditionedRestorer}) around PromptIR provides per-task affine modulation. Two learnable embeddings (one for rain, one for snow, each with scale and bias for 3 input channels) modulate the input before PromptIR and add a residual bias after. The wrapper adds only 18 parameters (2 embeddings $\times$ 3 channels $\times$ 3 tensors) while allowing the model to specialize its behavior per degradation type. At test time, since degradation labels are unknown, inference supports \texttt{--task\_mode both}, running both conditioned paths and averaging the outputs.

\subsubsection{SIPL-Lite Two-Pass Refinement}
Inspired by SIPL (Self-Improved Privilege Learning) \cite{wu2025}, a lightweight two-pass training scheme was implemented. After the first restoration pass, the second pass input blends the first-pass output with the clean image during training: $\text{input}_2 = (1-\alpha) \cdot \text{restored}_1 + \alpha \cdot \text{clean}$, where $\alpha$ linearly decays from a start value to zero. Total loss is $\mathcal{L} = \mathcal{L}(\text{restored}_1, \text{clean}) + \gamma \cdot \mathcal{L}(\text{restored}_2, \text{clean})$. At inference, the second pass uses the clamped first-pass output as input, enabling iterative self-refinement without ground-truth access.

\subsubsection{Per-Task Validation Metrics}
The validation dataset was extended to return degradation IDs ($0$ for snow, $1$ for rain), enabling separate logging of \texttt{val\_psnr\_derain}, \texttt{val\_psnr\_desnow}, \texttt{val\_loss\_derain}, and \texttt{val\_loss\_desnow} alongside the aggregate metrics. This provides visibility into whether experimental changes benefit one degradation type at the expense of another.

\subsubsection{Time-Estimate Progress Bar}
A custom Lightning \texttt{TQDMProgressBar} callback was added, printing an epoch-end timestamp line showing total elapsed wall-clock time and estimated remaining training time. This provides a clean, terminal-width-independent summary of training progress.

% ==========================================
% 4. RESULTS AND DISCUSSION
% ==========================================
\section{Results and Discussion}

\subsection{Experiment Progression and Leaderboard Scores}
Table \ref{tab:results} summarizes all experiments and their CodaBench public PSNR scores. Starting from an initial baseline of 30.52, systematic improvements raised the best score to 31.77 — a gain of 1.25 dB.

\begin{table}[htbp]
\centering
\caption{Progression of experiments and CodaBench public PSNR scores. Bold indicates current best.}
\label{tab:results}
\small
\setlength{\tabcolsep}{3pt}
\begin{tabular}{@{}clc>{\raggedright\arraybackslash}p{4.8cm}c@{}}
\toprule
\textbf{Run} & \textbf{Config} & \textbf{Ep.} & \textbf{Inference / Notes} & \textbf{PSNR} \\ \midrule
1 & p128, L1, bs=8 & 149 & original & 30.52 \\
2 & 5090 crash, early ckpt  & 09  & original & 24.93 \\
3 & p128, current scripts & 149 & original & 30.23 \\
4 & p192, bs=2 & 144 & original / TTA & 30.93 / 31.35 \\
5 & p192, bs=2, 200 ep & 144 & TTA (no gain over Run 4) & 31.35 \\
6 & p256, bs=1, L1 & 124 & original / TTA & 30.79 / 31.64 \\
7 & p384, bs=1, EMA & 109 & original / TTA / avg3 TTA & 30.78 / 31.49 / 31.50 \\
8 & p256, bs=1, \texttt{val\_psnr} & 149 & original / TTA / avg3 TTA & 31.27 / 31.73 / 31.73 \\
\textbf{9} & \textbf{p256, bs=1, L1+MSE ($\lambda$=0.05)} & \textbf{139} & \textbf{original / TTA / avg5 TTA} & \textbf{31.31 / 31.77 / 31.76} \\
10 & p256, L1+MSE + task cond. & 139 & original / TTA / avg5 TTA & 30.73 / 31.19 / 31.18 \\
11 & p256, L1+MSE + EMA & 129 & original / TTA / avg5 TTA & 31.28 / 31.72 / 31.73 \\
12 & p256, L1+MSE ($\lambda$=0.10) & — & running; pending submission & TBD \\ \bottomrule
\end{tabular}
\end{table}

\subsection{Key Findings}

\subsubsection{Patch Size Optimization (Runs 1–8)}
Early experiments explored patch sizes of 128, 192, 256, and 384. Since all images are $256 \times 256$, patch sizes larger than 256 are equivalent to full-image training; they do not provide additional spatial context. The $P=256$ configuration consistently outperformed $P=192$, likely because full-image context is more valuable for restoration than the higher data diversity from smaller random crops. Run 8 (PSNR-monitored checkpointing at $P=256$) achieved 31.73 with TTA and tied that score with top-3 checkpoint averaging.

\subsubsection{Loss Function: L1 vs L1+MSE (Run 9)}
The most impactful change was shifting from pure L1 loss to a hybrid L1+MSE loss with $\lambda=0.05$. Since the evaluation metric is PSNR, which directly depends on pixel-wise MSE, adding a small MSE term to the training objective better aligns the optimization target with the evaluation metric. Run 9 improved public TTA PSNR from 31.73 to 31.77 (+0.04 dB), and the original (no-TTA) score from 31.27 to 31.31.

\subsubsection{Task Conditioning Degrades Performance (Run 10)}
Run 10 added the \texttt{TaskConditionedRestorer} wrapper with L1+MSE loss. Despite the hypothesis that per-task modulation would help the model handle rain and snow differently, public PSNR collapsed to 31.19 (TTA) — a 0.58 dB drop. Local validation per-task metrics from Run 10 showed both derain and desnow PSNR were lower than Run 9, confirming that the conditioning hurt all degradation types. The likely cause is that the small embedding-based modulation overfits the known train/val labels and fails at test time, where forced task-mode averaging produces poor results because the conditioned model relies too heavily on the task label.

\subsubsection{EMA Does Not Help (Run 11)}
Exponential Moving Average (EMA) weight averaging was tested in isolation on the best Run 9 recipe. Run 11 with EMA achieved 31.72 (TTA), below the 31.77 baseline. EMA was previously tested on Run 7 ($P=384$) where it also did not improve over single-checkpoint TTA. The consistent negative result suggests that the PromptIR training dynamics do not benefit from EMA-based smoothing at this dataset scale.

\subsubsection{Checkpoint Averaging: Mixed Results}
Top-3 checkpoint averaging (avg3 TTA) tied single-checkpoint TTA at 31.73 in Run 8. Top-5 averaging (avg5 TTA) in Run 9 scored 31.76, slightly below the single-checkpoint TTA of 31.77. This suggests that the best single checkpoint captures near-optimal performance and averaging with nearby checkpoints does not provide additional robustness in the PromptIR loss landscape at 150 epochs.

\subsubsection{Effectiveness of TTA}
8-way geometric TTA consistently provided a substantial boost: $+0.48$ dB on average across all runs. This is significantly larger than typical TTA gains in classification, reflecting the importance of geometric consistency in restoration networks trained with random D4 augmentations. Run 9 TTA (31.77) vs original (31.31) represents a +0.46 dB gain.

\subsubsection{Validation vs Public Score Correlation}
Table \ref{tab:val_vs_public} compares best local validation PSNR with CodaBench public scores. Local validation PSNR is computed on the 320-image validation set, while public PSNR is computed on the 100-image hidden test set. The correlation is imperfect but informative.

\begin{table}[htbp]
\centering
\caption{Local validation PSNR vs CodaBench public PSNR. Higher local validation generally correlates with higher public scores, but the relationship is not monotonic.}
\label{tab:val_vs_public}
\small
\begin{tabular}{@{}lccc@{}}
\toprule
\textbf{Run} & \textbf{Local val\_psnr} & \textbf{Public TTA PSNR} & \textbf{Val–Public Gap} \\ \midrule
8  (L1, p256)     & 30.7184 & 31.73 & 1.01 \\
9  (L1+MSE $\lambda$=0.05) & 30.7490 & 31.77 & 1.02 \\
10 (L1+MSE + task)          & 30.7306 & 31.19 & 0.46 \\
11 (L1+MSE + EMA)           & 30.7558 & 31.72 & 0.96 \\ \bottomrule
\end{tabular}
\end{table}

The Run 10 case is notable: local validation PSNR was only 0.018 dB below Run 9, but public PSNR dropped by 0.58 dB — a 6$\times$ larger gap. This confirms that task conditioning overfits the visible validation split, appearing reasonable locally but collapsing on hidden test data.

\subsubsection{Per-Task Validation Analysis}
Runs 9–11 provide per-task validation metrics (introduced as a measurement-first modification). Table \ref{tab:pertask} shows the per-task breakdown from Run 9.

\begin{table}[htbp]
\centering
\caption{Per-task validation metrics from Run 9 (L1+MSE, best configuration). Desnow consistently outperforms derain by over 2 dB.}
\label{tab:pertask}
\small
\begin{tabular}{@{}lccc@{}}
\toprule
\textbf{Metric} & \textbf{Aggregate} & \textbf{Derain} & \textbf{Desnow} \\ \midrule
val\_psnr  & 30.749 & 29.654 & 31.887 \\
val\_loss  & 0.0196 & 0.0236 & 0.0156 \\ \bottomrule
\end{tabular}
\end{table}

Derain is the performance bottleneck: desnow achieves nearly 32 dB while derain struggles at 29.7 dB. This 2.2 dB gap suggests that the current PromptIR training strategy is better suited for snow removal than rain streak removal. Future work should focus specifically on improving derain performance, potentially through stronger degradation-specific augmentation or modified loss weighting for rain images.

\subsubsection{Failed Approaches}
Several strategies were evaluated and found to be harmful or ineffective:
\begin{itemize}
    \item \textbf{Stage 2 Merged Retraining:} Training from scratch on Train+Val merged data with no validation monitor degraded PSNR from 30.5 to 21.7 (Runs 1–3). The loss of a validation checkpoint selection mechanism makes this approach fundamentally unreliable on small datasets.
    \item \textbf{Task Conditioning (Run 10):} Despite the theoretical appeal of per-task specialization, the simple embedding-based approach overfits and transfers poorly to hidden test labels.
    \item \textbf{EMA (Run 11):} Two independent tests (Runs 7 and 11) showed EMA degrades or neutralizes scores, indicating PromptIR training does not benefit from weight averaging.
    \item \textbf{SIPL-Lite Inference:} A second-pass refinement during inference was tested on Run 11 and scored 31.72, below the single-pass TTA of 31.77.
    \item \textbf{Top-N Checkpoint Averaging:} Both avg3 and avg5 tied or slightly underperformed the best single checkpoint, suggesting the PromptIR loss landscape does not contain complementary checkpoints near the optimum.
\end{itemize}

% ==========================================
% 5. CONCLUSION
% ==========================================
\section{Conclusion}

This project successfully trained a PromptIR model from scratch for all-in-one rain and snow image restoration. Starting from a baseline of 30.52 PSNR, systematic experimentation yielded a final best score of \textbf{31.77 PSNR} (L1+MSE loss, patch 256, TTA), representing a 1.25 dB improvement.

The key contributions are:
\begin{enumerate}
    \item \textbf{Loss Alignment:} Replacing pure L1 loss with L1+$\lambda$MSE ($\lambda$=0.05) aligned the training objective with the PSNR evaluation metric, yielding +0.04 dB public PSNR (Run 8 $\rightarrow$ Run 9).
    \item \textbf{PSNR-Monitored Checkpointing:} Switching from \texttt{val\_loss} to \texttt{val\_psnr} monitoring improved checkpoint selection, raising original PSNR from 30.79 to 31.27 (+0.48 dB).
    \item \textbf{Geometric TTA:} 8-way flip/transpose test-time augmentation consistently provided +0.46 dB over single-original inference, accounting for the largest individual gain in the submission pipeline.
    \item \textbf{Patch Size Tuning:} $P=256$ full-image training outperformed $P=192$ and $P=384$ configurations by up to 0.38 dB.
    \item \textbf{Per-Task Diagnostics:} Separate derain/snow validation metrics revealed derain as the performance bottleneck (2.2 dB below desnow), guiding future research direction.
    \item \textbf{Identification of Failed Strategies:} Task conditioning, EMA, SIPL-lite inference, and Stage 2 merged retraining were all systematically tested and found to degrade or provide no improvement over the baseline.
\end{enumerate}

The fact that task conditioning and EMA both failed despite strong theoretical motivation highlights an important lesson: on small paired datasets trained from scratch, simple recipe improvements (loss function, checkpoint selection, TTA) often outperform architectural complications. The PromptIR transformer backbone is already a strong restoration model; squeezing additional gains requires careful empirical tuning rather than module additions that risk overfitting.

\subsubsection{Recommendations for Further Improvement}
Given the remaining 2.2 dB gap between derain and desnow, future experiments should focus on degradation-specific enhancements:
\begin{itemize}
    \item Stronger MSE weighting ($\lambda = 0.10$ or $0.20$) to better align with PSNR.
    \item Rain-specific data augmentation (rain-streak augmentation, intensity variation).
    \item Longer training (200+ epochs) with aggressive cosine restarts.
    \item Mixed-precision training (FP16) to enable higher batch sizes.
\end{itemize}

The submission with the best public PSNR is \texttt{submission/stage1-p256-tta-run9-l1mse.zip} at 31.77.

% ==========================================
% REFERENCES
% ==========================================
\newpage
\begin{thebibliography}{99}
\bibitem{potlapalli2023} Potlapalli, V., Zamir, S.W., Khan, S., Khan, F.S.: PromptIR: Prompting for All-in-One Blind Image Restoration. In: Proc. NeurIPS (2023)
\bibitem{valanarasu2022} Valanarasu, J.M.J., Yasarla, R., Patel, V.M.: TransWeather: Transformer-based Restoration of Images Degraded by Adverse Weather Conditions. In: Proc. CVPR (2022)
\bibitem{zamir2022} Zamir, S.W., Arora, A., Khan, S., Hayat, M., Khan, F.S., Yang, M.H.: Restormer: Efficient Transformer for High-Resolution Image Restoration. In: Proc. CVPR (2022)
\bibitem{li2022} Li, B., Liu, X., Hu, P., Wu, Z., Lv, J., Peng, X.: All-in-One Image Restoration for Unknown Corruption. In: Proc. CVPR (2022)
\bibitem{ashukha2020} Ashukha, A., Lyzhov, A., Molchanov, D., Vetrov, D.: Pitfalls of In-Domain Uncertainty Estimation and Ensembling in Deep Learning. In: Proc. ICLR (2020)
\bibitem{horton2021} Wortsman, M., Ilharco, G., Gadre, S.Y., Roelofs, R., et al.: Model soups: averaging weights of multiple fine-tuned models improves accuracy without increasing inference time. In: Proc. ICML (2022)
\bibitem{wu2025} Wu, G., Jiang, J., Jiang, K., Liu, X.: Boosting All-in-One Image Restoration via Self-Improved Privilege Learning. arXiv:2505.24207 (2025)
\end{thebibliography}

\end{document}
